% Brief report, 2026-08-17. Compile from the PROJECT ROOT (figure paths are root-relative):
%   xelatex -output-directory=presentation presentation/08_17.tex   (run twice)
\documentclass[11pt]{article}
\usepackage[margin=1in]{geometry}
\usepackage{amsmath}
\usepackage{graphicx}
\usepackage[hidelinks]{hyperref}
\usepackage[font=small,labelfont=bf]{caption}

\begin{document}

\section{Capped vs.\ uncapped mean earnings}

Extrapolating the fitted cross-sections forward in time, we hold the \emph{shape} of each
cell's distribution constant and shift only its location with the wage index. To see
whether that is a reasonable assumption we compare the ASS uncapped and capped (taxable)
mean earnings per worker: with the taxable maximum indexed to wages (1975 onward), a
constant shape keeps this ratio flat, so its drift reflects shape change --- in
particular in the upper tail. Figure~\ref{fig:ratio} shows the ratio over 1937--2022.

\begin{figure}[h!]\centering
\includegraphics[width=.8\textwidth]{output/cross_sections/plots/ass_uncapped_capped_ratio.pdf}
\caption{Ratio of ASS uncapped to capped (taxable) mean covered earnings per worker.}
\label{fig:ratio}
\end{figure}

\section{Estimation}

Each (year, sex, single-year age) cell's objective is a \emph{convex combination} of the
censored log-likelihood of the EPUF microdata and a GMM criterion on the corresponding
cohort~$\times$~age statistics of Guvenen, Kaplan, Song and Weidner (AEJ:\ Applied 2022;
GKSW), with weight $\lambda=\tfrac12$ on each part. The GMM side targets ten statistics
per cell: the mean, standard deviation, skewness and kurtosis of log earnings, and the
$p10$, $p25$, $p50$, $p75$, $p90$, $p98$ quantiles, all for earnings above GKSW's
selection floor of $260\times$ the nominal minimum wage. The fitted distributions are a
double Pareto-lognormal for men and a two-component lognormal mixture for women.

\paragraph{The two datasets differ.} GKSW's statistics come from the 1\% CWHS and use
wage and salary income (W-2 Box~1 from 1978) of commerce-and-industry workers ---
self-employment is excluded (their footnote~10: available only from 1978 and top-coded
until 1994), and above-cap earnings are uncensored from 1978 (imputed before) ---
winsorized at the 99.999th percentile. EPUF records Social-Security-covered earnings of
\emph{all} covered workers, including covered self-employment, top-coded at the taxable
maximum and random-rounded. Figure~\ref{fig:qval} computes the same quantiles from both
sources under the identical selection floor, aggregated by cohort over the complete age
range 25--55 and using only cohorts with no cell censored at the cap; $p98$ (and $p90$
for men) are left out because they are never observed in EPUF.

\begin{figure}[h!]\centering
\includegraphics[width=.92\textwidth]{output/cross_sections/plots/guv_quantile_validation_men.pdf}\\[2pt]
\includegraphics[width=.92\textwidth]{output/cross_sections/plots/guv_quantile_validation_women.pdf}
\caption{GKSW sel0 quantiles vs the same quantiles in raw EPUF (men top, women bottom).}
\label{fig:qval}
\end{figure}

\section{Model vs.\ GKSW (2022)}

Figure~\ref{fig:guvcomp} aggregates the ten targeted statistics by cohort --- an
equal-weight average over each cohort's ages 25--55 --- and plots the GKSW data against
the same functionals computed from the fitted model, for men and women.

\begin{figure}[h!]\centering
\includegraphics[width=.95\textwidth]{output/cross_sections/plots/guv_comparison_men_guvgmm_smoothed.pdf}\\[2pt]
\includegraphics[width=.95\textwidth]{output/cross_sections/plots/guv_comparison_women_guvgmm_smoothed.pdf}
\caption{GKSW cohort~$\times$~age functionals (blue) vs the fitted model (orange), by
cohort: men (top) and women (bottom).}
\label{fig:guvcomp}
\end{figure}

\section{Aggregates}

Figure~\ref{fig:agg} shows total taxable (capped) earnings and the uncapped mean
earnings per worker implied by the model against the published ASS series (spliced onto
the Trustees Report projection after 2022). Neither moment is targeted by the
estimation. Cells are aggregated with the observed EPUF (sex, age) composition of
positive earners in each year, scaled to the published covered-worker totals; off-sample
years hold the composition fixed at the nearest observed edge (the 1951--55 average
backward, the 2000--04 average forward).

\begin{figure}[h!]\centering
\includegraphics[width=.98\textwidth]{output/cross_sections/plots/aggregate_taxable_extrapolated_guvgmm_extrapolated.pdf}
\caption{Aggregate taxable earnings vs ASS+TR (top row) and uncapped mean earnings per
worker vs ASS (bottom row), 1937--2100.}
\label{fig:agg}
\end{figure}

\end{document}
